\section{Main Construction}
\label{sec:constr}

\subsection{Construction of Circle Bundles with \psc}

In the beginning we raised the following question.

\begin{question}
	Does there exist a circle bundle $E \to M$ over an enlargeable manifold $M$, such that the total space $E$ admits a \psc-metric?
\end{question}

We will now construct examples of such bundles in all dimensions $\mathrm{dim}(M) \geq 4$. Let us first describe the strategy behind the construction. The existence of a \psc-metric on $E$ will be established as an application of Theorem \ref{psc:thm:existence_result_totally_non_spin}, restated here for convenience:

\vspace{\baselineskip}
\noindent\textbf{Theorem~\ref{psc:thm:existence_result_totally_non_spin}.}\itshape\ 
Let $W^{n+1}$ be an oriented cobordism from $M_0$ to $M_1$, $n \geq 5$. Suppose that $M_1$ is a connected totally non-spin manifold, the inlcusion $M_1 \to W$ induces an isomorphism on $\pi_1$ and that $M_0$ carries a psc-metric. Then $M_1$ also carries a psc-metric. In fact the psc-metric on $M_0$ can be extended to a psc-metric on $W$, which is of product form near the boundary. 
\normalfont
\vspace{\baselineskip}

In order to apply this in our case, we make the following two observations, which hold true for any circle bundle $E \to M$ over a closed orientable manifold $M$. 

\begin{enumerate}[label=(\roman*)]
	\item There exists an oriented nullbordism of $E$, namely the total space of the associated disk bundle $DE \to M$,
	\item \label{constr:eq:circle_bundle_observation2} The inclusion $E \cong \partial DE \to DE$ induces an isomorphism on $\pi_1$ if and only if the fibers of $E$ are contractible, that is the inclusion $S^1 \to E$ is trivial on $\pi_1$.
\end{enumerate}

We have already observed in Example \ref{bundles:ex:associated_bundles_part3} that $DE$ is a nullbordism of $E$. Moreover, the structure group of $DE \to M$ acts on fibers by orientation preserving diffeomorphisms. Thus any orientation on $M$ gives rise to an orientation on $DE$, which is uniquely determined by requiring local trivialisations to be orientation preserving. For \ref{constr:eq:circle_bundle_observation2} note that $DE \to M$ is a homotopy equivalence. Hence the diagram 
\[\begin{tikzcd}[column sep=0cm]
	E && DE \\
	& M
	\arrow[from=1-1, to=1-3]
	\arrow[from=1-1, to=2-2]
	\arrow["\simeq", from=1-3, to=2-2]
\end{tikzcd}\]
shows that $E \to DE$ induces an isomorphism on $\pi_1$ if and only if $E \to M$ does. Then the claim follows follows from the long exact sequence of homotopy groups:
\[\begin{tikzcd}
	{\pi_1(S^1)} & {\pi_1(E)} & {\pi_1(M)} & {\pi_0(S^1)}
	\arrow[from=1-1, to=1-2]
	\arrow[from=1-2, to=1-3]
	\arrow[from=1-3, to=1-4]
\end{tikzcd}\]

So to answer the question above positively, we have to construct circle bundles $E \to M$ with the following properties.

\begin{enumerate}[label = (\Roman*)]
	\item \label{constr:property1} $M$ is enlargeable, orientable and connected,
	\item \label{constr:property2} $E$ is totally non-spin,
	\item \label{constr:property3} The inclusion $S^1 \to E$ is trivial on $\pi_1$.
\end{enumerate}

We define $M$ as the connected sum of three closed manifolds. Let $N^{2n}$ be any $2n$-dimensional orientable totally non-spin manifold. For example take $N:= \CP^2 \times T^{2n-4}$, which is totally non-spin by Lemma \ref{psc:lem:totally_non_spin_lemma} and Example \ref{bundles:ex:cp2_not_spin}. Then consider
\[
M := \CP^n \# T^{2n} \# N.
\] 
Intuitively the bundle $E \to M$ will now look as follows. It is trivial over the summands $T^{2n}$ and $N$, and is given by the canonical bundle $S^{2n+1} \to \CP^n$ over the $\CP^n$ summand (or rather the restriction of this bundle to $\CP^n - D^{2n} \subseteq M$). Schematically:

\placeholder

The presence of the Torus summand then corresponds to property \ref{constr:property1}, the trivial bundle over $N$ to property \ref{constr:property2} and the $S^{2n+1}$ over $\CP^n$ to property \ref{constr:property3}. We will now construct this bundle rigourously and afterwards discuss the proof of properties \ref{constr:property1} - \ref{constr:property3} in more detail. 

Let 
\[
c_1 \in H^2(M;\Z) \cong H^2(\CP^n;\Z) \oplus H^2(T^{2n};\Z) \oplus H^2(N;\Z)
\]
be the cohomology class corresponding to the first chern class $\alpha \in H^2(\CP^n;\Z)$ of the canonical bundle $S^{2n+1} \to \CP^n$ under the above isomorphism. Define $E \to M$ as the unique $S^1$-principal bundle with first chern class equal to $c_1$. We claim that $E$ is as described above.



Since the collapse map $M \to T^{2n}$ has degree $1$ and $T^{2n}$ is enlargeable, it follows from Lemma \ref{enl:lem:constructing_enlargeable_mnf} that $M$ is also enlargeable. 



For $n \geq 2$  Now consider the manifold

Since the collapse map $M \to T^{2n}$ has degree $1$ and $T^{2n}$ is enlargeable, it follows from Lemma \ref{enl:lem:constructing_enlargeable_mnf} that $M$ is also enlargeable. We will now construct a circle bundle $E \to M$ as follows. Let 
\[
c_1 \in H^2(M;\Z) \cong H^2(\CP^n;\Z) \oplus H^2(T^{2n};\Z) \oplus H^2(N;\Z)
\]
be the cohomology class corresponding to the first chern class $\alpha \in H^2(\CP^n;\Z)$ of the canonical bundle $S^{2n+1} \to \CP^n$ under the above isomorphism. We then define $E \to M$ as the unique $S^1$-principal bundle with first chern class equal to $c_1$. Now this is quite an abstract definition, so let us develop an intuition for how this bundle looks like. We show that $E$ is trivial over the summand $T^{2n} \# N$ and is given by the restriction of the canonical bundle $S^{2n+1} \to \CP^n$ to $\CP^n - D^{2n} \subseteq M$ over the $\CP^n$ summand. To make this precise let 
\[
i_1 \colon (T^{2n} \# N) - D^{2n} \to M
\]
be the inclusion. Then
\[\begin{tikzcd}
	{H^2(M; \Z)} & {H^2((T^{2n} \# N) - D^{2n};\Z) } \\
	{H^2(\CP^n;\Z) \oplus H^2(T^{2n} \# N;\Z)} & {H^2(T^{2n} \# N;\Z) }
	\arrow["{i_1^\ast}", from=1-1, to=1-2]
	\arrow["\cong", from=2-1, to=1-1]
	\arrow["{\mathrm{proj.}}", from=2-1, to=2-2]
	\arrow["\cong", from=2-2, to=1-2]
\end{tikzcd}\]
commutes and therefore $i_1^\ast (c_1) = 0$. Again using the fact that $S^1$-principal bundles are uniquely determined by their first chern class we see that $i_1^\ast E$ is indeed trivial. Similiarly if we denote by 
\[
i_2 \colon \CP^n - D^{2n} \to M, \quad j \colon \CP^{n}-D^{2n} \to \CP^{n}
\]
the inclusions, the diagram
\[\begin{tikzcd}
	{H^2(M ;\Z)} & {H^2(\CP^n - D^{2n};\Z) } \\
	{H^2(\CP^n;\Z) \oplus H^2(T^{2n} \# N;\Z)} & {H^2(\CP^n;\Z) }
	\arrow["{i_2^\ast}", from=1-1, to=1-2]
	\arrow["\cong", from=2-1, to=1-1]
	\arrow["{\mathrm{proj.}}", from=2-1, to=2-2]
	\arrow["\cong", "j^\ast"', from=2-2, to=1-2]
\end{tikzcd}\]
commutes and thus $i_2^\ast(c_1) = j^\ast(\alpha)$. So the first chern classes of $i_2^\ast E$ and $j^\ast S^{2n+1}$ agree, hence they must be isomorphic bundles. 

With this intuition in mind we will now show the following two claims.

\begin{claim}
	\label{claim1}
$E$ is totally non-spin.
\end{claim}

\begin{claim}
	\label{claim2}
The projection $p \colon E \to M$ induces an isomorphism on $\pi_1$.
\end{claim}

Granting these claims for now, let us see how to conclude the existence of a positive scalar curvature metric on $E$. The disk bundle $DE \to M$ corresponding to $E$ defines an oriented null-cobordism of $E$ (see Example \ref{bundles:ex:associated_bundles} \ref{bundles:ex:associated_bundles_part3}). Moreover, since $DE \to M$ is a homotopy equivalence it follows from Claim \ref{claim1} that the inclusion $E \to DE$ is an isomorphism on $\pi_1$. So together with Claim \ref{claim2} Theorem \ref{psc:thm:existence_result_totally_non_spin} applies, hence $E$ admits a psc metric.

Let us now proof Claim \ref{claim1} and Claim \ref{claim2}.

% Granting these Claims for now let us see why they suffice for the assumptions of Theorem \ref{thm:gromov_hanke} to be satisfied.  
% For this we would like to show that the classifying map $\phi \colon E \to B\pi_1(E)$ maps $[E]$ to $0 \in H_{2n+1}(\pi_1(E);\Z)$, which can be seen by using Claim \ref{claim2} as follows. Let $q \colon DE \to M$ be the disk bundle corresponding to $E$. Moreover let $\Phi \colon DE \to B\pi_1(DE)$ be the classifiying map for $DE$. Then the diagram 
% \[\begin{tikzcd}
% 	DE & {B\pi_1(DE)} \\
% 	E & {B\pi_1(E)}
% 	\arrow["\Phi", from=1-1, to=1-2]
% 	\arrow["\iota", from=2-1, to=1-1]
% 	\arrow["\phi", from=2-1, to=2-2]
% 	\arrow["{B\pi_1(\iota)}"', from=2-2, to=1-2]
% \end{tikzcd}\]
% commutes (hier noch Begründung rein), where the vertical maps are (induced by) the inclusion $\iota \colon E \to DE$. We claim that $B\pi_1(\iota)$ is an isomorphism. This follows from commutativity of  
% \[\begin{tikzcd}[column sep = small]
% 	{B\pi_1(E)} && {B\pi_1(DE)} \\
% 	& {B\pi_1(M)}
% 	\arrow["{B\pi_1(\iota)}", from=1-1, to=1-3]
% 	\arrow["{B\pi_1(p)}"',"\cong", from=1-1, to=2-2]
% 	\arrow["{B\pi_1(q)}","\cong"', from=1-3, to=2-2]
% \end{tikzcd}\]
% where $B\pi_1(q)$ is an isomorphism by homotopy invariance and $B\pi_1(p)$ is an isomorphism by Claim \ref{claim2}. Thus $\phi$ factors over $\iota$. Since $DE$ is an oriented null-cobordism of $E$ it follows that $\iota_\ast[E] = 0 \in H_{2n+1}(DE;\Z)$ and so in particular $\phi_\ast [E] = 0\in H_{2n+1}(\pi_1(E);\Z)$ as we wanted to show. So together with Claim \ref{claim1} all assumptions of Theorem \ref{thm:gromov_hanke} are satisfied, showing that $E$ admits a metric of positive scalar curvature. 

\begin{proof}[Proof of Claim \ref{claim1}]
	Let $i \colon N- D^{2n} \to M$ be the inclusion. We have already established above, that $E$ is trivial over the $T^{2n} \# N$ summand and so in particular
	\[
	i^\ast E \cong (N- D^{2n}) \times S^1.
	\]
	From Lemma \ref{lem:totally_non_spin_lemma} \ref{eq:totally_non_spin_lemma1} and \ref{eq:totally_non_spin_lemma3} we conclude that $(N- D^{2n}) \times S^1$ is totally non-spin. 
	% Then the universal cover of $i^\ast E$ is given by $\tilde{N}_0 \times \R$, where $\tilde{N}_0$ is the universal cover of $N_0$. Let $\pi_{\tilde{N}_0} \colon \tilde{N}_0 \times \R \to \tilde{N}_0$ and $\pi_\R \colon \tilde{N}_0 \times \R \to \R$ be the projections. We compute 
	% \[
	% w_2(T(\tilde{N}_0 \times \R)) = w_2(\pi_{\tilde{N}_0}^\ast T\tilde{N}_0 \oplus \pi_\R^\ast T\R) = \pi_{\tilde{N}_0}^\ast w_2(T\tilde{N}_0).
	% \]
	% Since $\pi_{\tilde{N}_0}$ is a homotopy equivalence and $\tilde{N}_0$ is non-spin it follows that $w_2(T(\tilde{N}_0 \times \R)) \neq 0$, such that $i^\ast E$ is totally non spin. 
	Therefore $E$ contains a totally non-spin submanifold of codimension $0$ and so Lemma \ref{lem:totally_non_spin_lemma} \ref{eq:totally_non_spin_lemma2} implies that $E$ is also totally non-spin.
\end{proof}

\begin{proof}[Proof of Claim \ref{claim2}]
	We have to show that the fibers of $E$ can be contracted. This is essentially due to the fact that $E$ contains (the restriction of) the canonical bundle $S^{2n+1} \to \CP^n$ for which the fibers are contractible. To formalize this let $j \colon \CP^{n} - D^{2n}$ be the inclusion. Consider the following commutative diagram, where the rows are taken from the long exact sequences of homotopy groups of the respective fibrations and the vertical maps are induced by inclusions: 
	\[\begin{tikzcd}
	{\pi_2(M)} & {\pi_1(S^1)} & {\pi_1(E)} \\
	{\pi_2(\CP^n - D^{2n})} & {\pi_1(S^1)} & {\pi_1(j^\ast E)} \\
	{\pi_2(\CP^n)} & {\pi_1(S^1)} & {\pi_1(S^{2n+1})}
	\arrow["{\partial_1}", from=1-1, to=1-2]
	\arrow[from=1-2, to=1-3]
	\arrow[from=2-1, to=1-1]
	\arrow["{\partial_2}", from=2-1, to=2-2]
	\arrow["\cong"', from=2-1, to=3-1]
	\arrow["{\mathrm{id}}", from=2-2, to=1-2]
	\arrow[from=2-2, to=2-3]
	\arrow["{\mathrm{id}}"', from=2-2, to=3-2]
	\arrow[from=2-3, to=1-3]
	\arrow[from=2-3, to=3-3]
	\arrow["{\partial_3}", from=3-1, to=3-2]
	\arrow[from=3-2, to=3-3]
\end{tikzcd}\]
The vertical map on the lower left is an isomorphism, because both $\CP^n$ and $\CP^n - D^{2n}$ are simply connected and so by the Hurewicz Theorem
\[\begin{tikzcd}
	{\pi_2(\CP^n-D^{2n})} & {\pi_2(\CP^n)} \\
	{H_2(\CP^n-D^{2n};\Z)} & {H_2(\CP^n;\Z)}
	\arrow[from=1-1, to=1-2]
	\arrow["\cong", from=1-1, to=2-1]
	\arrow["\cong", from=1-2, to=2-2]
	\arrow[from=2-1, to=2-2]
\end{tikzcd}\]
commutes and the vertical maps are isomorphisms. Since the bottom horizontal map is an isomorphism the same holds for the top horizontal map. We now claim that $\partial_1$ is surjective. For this it suffices to show surjectivity of $\partial_2$, which in turn is equivalent to surjectivity of $\partial_3$. But this is clear by exactness of the rows and $\pi_1(S^{2n+1}) = 0$. Consequently
\[
\pi_1(S^1) \to \pi_1(E)
\]
is the zero map and therefore we conclude from exactness of
\[
\pi_1(S^1) \to \pi_1(E) \to \pi_1(M) \to \pi_0(S^1),
\]
that $\pi_1(E) \to \pi_1(M)$ is indeed an isomorphism.
\end{proof}

The above construction now gives examples for circle bundles with positive scalar curvature over enlargeable manifolds $M$ in all even dimensions $\mathrm{dim}(M)\geq 4$. To also get examples in odd dimensions let $E \to M$ be as above and let $\pi \colon M \times S^1 \to M$ be the projection. Then the pullback bundle $\pi^\ast E \cong E \times S^1$ clearly admits a metric of positive scalar curvature and by Lemma \ref{enl:lem:constructing_enlargeable_mnf} the base space $M \times S^1$ is still enlargeable. Naturally, the question arises as to what happens in dimensions $\mathrm{dim}(M) \leq 3$. In this case it turns out that no circle bundle $E \to M$ admits a metric of positive scalar curvature, see \cite[Theorem B]{kumar2025circle}.


\subsection{Necessity of the Enlargeability Assumption in Gromov's Decay Theorem}


Coming back to..


\subsection{Further Questions}
